% concatenated from k-Lucas_Annulus.tex
\documentclass[11pt,a4paper]{article}
\usepackage[latin1]{inputenc}
\usepackage{amsmath}
\usepackage{amsthm}
\usepackage{amsfonts}
\usepackage{amsfonts,amsthm,latexsym,amsmath,amssymb,amscd,epsfig,psfrag,enumerate}
\usepackage{graphics,graphicx, bezier, float, color, hyperref}
\usepackage{amssymb,url}
\usepackage{multienum}
\usepackage[table]{xcolor}
\usepackage{multicol,multirow}
\usepackage{graphicx}
\usepackage{fancyvrb}
\usepackage{parskip}
\usepackage[toc,page]{appendix}
\sloppy
\setlength{\parindent}{0pt}
\setlength\parskip{0.1in}
\usepackage[top=2.7cm, bottom=2.7cm, left=1.5cm, right=1.5cm]{geometry}
\usepackage{xcolor}
\usepackage{blkarray}
\newtheorem{theorem}{Theorem}[section]
\newtheorem{lemma}{Lemma}[section]
\newtheorem{cor}{Corollary}[section]
\newtheorem{remark}{Remark}[section]
\usepackage[none]{hyphenat}[section]
\newtheorem{definition}{Definition}[section]
\newcommand\numberthis{\addtocounter{equation}{1}\tag{\theequation}}
\numberwithin{equation}{section}
\numberwithin{table}{section}
\numberwithin{figure}{section}
\title{$k$-Lucas Annulus for Polynomial Zeros}
\author{Herbert Batte$^{1,*}$}
\date{}
\begin{document}
	\maketitle
	\begin{center}
		\textit{In memory of Professor Florian Luca (1969--2026)}
	\end{center}
	\abstract{
We prove a binomial identity relating the $k$-Lucas and $k$-Fibonacci sequences: for real $k>0$ and integers $m\ge1$, $n\ge1$, a binomial-weighted sum of $k$-Lucas numbers reduces to a multiple of $L_{k,mn}$ when $n$ is even, but to a multiple of $F_{k,mn}$ when $n$ is odd. This parity dependence traces back to the $k$-Lucas Binet formula lacking the normalising factor $1/(\alpha-\beta)$ present for $k$-Fibonacci numbers. The identity supplies, for each parity of $n$ and each $m$, an explicit family of positive weights summing to one; combined with the general annulus principle of Dalal and Govil, this yields a $k$-Lucas annulus containing all the zeros of a complex polynomial of degree $n$, complementing the $k$-Fibonacci-based bounds of D\'iaz-Barrero, Bidkham-Shashahani, and Kaur.
	}
	
	{\bf Keywords and phrases}: $k$-Lucas numbers; $k$-Fibonacci numbers; location of zeros; annulus; binomial coefficients.
	
	{\bf 2020 Mathematics Subject Classification}: 30C15, 30C10, 11B39.
	
	\thanks{$ ^{*} $ Corresponding author}
	
	\section{Introduction}
	
\subsection{Background}\label{sec:background}
Cauchy's classical theorem places all zeros of a polynomial
\begin{align*}
q(z)=\sum_{j=0}^{n}a_jz^j,\qquad a_n\neq 0,	
\end{align*}
 inside the disc 
 \begin{align*}
|z|\le 1+\max_{0\le j<n}|a_j/a_n|. 	
 \end{align*}
 Since then, sharper localisations have been obtained by weighting the coefficients of $q$ against an auxiliary sequence and producing an \emph{annulus} $r_1\le|z|\le r_2$ rather than a single disc. D\'iaz-Barrero \cite{diazbarrero} obtained such an annulus built from the ordinary Fibonacci numbers and binomial coefficients. Bidkham and Shashahani \cite{bidkham} generalised this by replacing the Fibonacci numbers with the one-parameter $k$-Fibonacci sequence, producing a family of annuli indexed by $k>0$ and recovering D\'iaz-Barrero's bound at $k=1$. Kaur \cite{kaur} generalised the same result along a different axis, replacing the fixed exponent $4$ that appears in both earlier bounds by an arbitrary positive integer $m$, and recovering the Bidkham--Shashahani annulus as the special case $m=4$. Underlying all three constructions is the same mechanism: given any positive weights $B_1,\dots,B_n$ summing to $1$, a reverse-triangle-inequality argument applied to $q(z)$ and to $z^nq(1/z)$ produces an annulus from those weights alone. Dalal and Govil \cite{dalalgovil} isolated this principle explicitly, so that the substance of a result in this line lies not in the annulus-construction mechanism itself but in the specific closed-form weights it supplies.

All three results rest on the $k$-Fibonacci sequence of Falc\'on and Plaza \cite{falconplaza}, defined by
 \begin{align*}
F_{k,0}=0,\qquad F_{k,1}=1,\qquad F_{k,n+1}=kF_{k,n}+F_{k,n-1}\quad(n\ge1),
 \end{align*}
where $k$ is a fixed positive real number, which admits the Binet form
 \begin{align*}
F_{k,n}=\frac{\alpha^n-\beta^n}{\alpha-\beta},\qquad
\alpha=\frac{k+\sqrt{k^2+4}}{2},\quad \beta=\frac{k-\sqrt{k^2+4}}{2},
 \end{align*}
with $\alpha,\beta$ the two roots of $x^2-kx-1=0$. The companion sequence sharing the same characteristic roots is the $k$-Lucas sequence introduced by Falc\'on \cite{falconlucas},
 \begin{align*}
L_{k,0}=2,\qquad L_{k,1}=k,\qquad L_{k,n+1}=kL_{k,n}+L_{k,n-1}\quad(n\ge1),
 \end{align*}
with Binet form $L_{k,n}=\alpha^n+\beta^n$ and the same $\alpha,\beta$ as above, but \emph{without} the normalising factor $1/(\alpha-\beta)$ that the $k$-Fibonacci sequence carries. This missing factor turns out not to be cosmetic. 

The identity underlying the D\'iaz-Barrero, Bidkham-Shashahani and Kaur annulus is, at heart, a binomial expansion of $(F_{k,m-1}+\alpha F_{k,m})^n$ that collapses to a single power of $\alpha$ precisely because the $1/(\alpha-\beta)$ factors carried by $F_{k,m-1}$ and $F_{k,m}$ cancel. Repeating the same computation with $L_{k,m-1}$ and $L_{k,m}$ in place of $F_{k,m-1}$ and $F_{k,m}$, no such cancellation is available, and the resulting sum depends on the parity of $n$. This reduces to a multiple of $L_{k,mn}$ when $n$ is even, and to a multiple of $F_{k,mn}$ (not $L_{k,mn}$), when $n$ is odd. We record this precisely as Theorem \ref{thm:main} below, and use it to supply, in each of the two parity cases, an explicit new choice of weights for the Dalal-Govil principle, yielding a $k$-Lucas annulus containing all the zeros of $q(z)$.

	
\subsection{Main Results}

\begin{theorem}[$k$-Lucas -- $k$-Fibonacci binomial identity]
	\label{thm:main}
	Let $k>0$ be a real number, and let $m,n$ be integers with $m\ge 1$ and $n\ge1$. Then
 \begin{align*}
	\sum_{\ell=0}^{n}\binom{n}{\ell}(L_{k,m-1})^{n-\ell}(L_{k,m})^{\ell}L_{k,\ell}
	=
	\begin{cases}
		(k^2+4)^{n/2}\,L_{k,mn}, & \text{if } n \text{ is even},\\[6pt]
		(k^2+4)^{(n+1)/2}\,F_{k,mn}, & \text{if } n \text{ is odd}.
	\end{cases}
 \end{align*}
\end{theorem}

\begin{remark}
	The parity dependence in Theorem \ref{thm:main} is not an artefact of the proof. It reflects the same mechanism behind the relation $L_{k,n}^2=(k^2+4)F_{k,n}^2+4(-1)^n$ established by Falc\'on \cite{falconlucas}, in which the sign separating $L_{k,n}^2$ from $(k^2+4)F_{k,n}^2$ already depends on the parity of $n$. Theorem \ref{thm:main} shows the same phenomenon propagates into the binomial convolution that, for $k$-Fibonacci numbers, is parity-independent.
\end{remark}
	
\begin{definition}
	\label{def:Dn}
	For integers $m\ge1,n\ge1$ and real $k>0$, set
 \begin{align*}
	D_n=D_n(k,m):=
	\begin{cases}
		(k^2+4)^{n/2}L_{k,mn}-2(L_{k,m-1})^n, & \text{if } n \text{ is even},\\[6pt]
		(k^2+4)^{(n+1)/2}F_{k,mn}-2(L_{k,m-1})^n, & \text{if } n \text{ is odd}.
	\end{cases}
 \end{align*}
\end{definition}

By Theorem \ref{thm:main}, $D_n=\sum_{\ell=1}^{n}\binom{n}{\ell}(L_{k,m-1})^{n-\ell}(L_{k,m})^{\ell}L_{k,\ell}$, a sum of strictly positive terms; in particular $D_n>0$ always (proved formally in Section \ref{prel}).

\begin{theorem}[$k$-Lucas annulus for polynomial zeros]
	\label{thm:annulus}
	Let $q(z)=a_0+a_1z+\cdots+a_nz^n$ be a polynomial of degree $n\ge1$ with complex coefficients, $a_\ell\neq0$ for $0\le\ell\le n$. Let $k>0$ be real and $m\ge1$ an integer. Then all zeros of $q$ lie in the annulus
 \begin{align*}
	D=\{z\in\mathbb{C}: r_1\le|z|\le r_2\},
	\end{align*}
	where
 \begin{align*}
	r_1=\min_{1\le \ell\le n}\left\{\binom{n}{\ell}\frac{(L_{k,m-1})^{n-\ell}(L_{k,m})^{\ell}L_{k,\ell}}{D_n}\left|\frac{a_0}{a_\ell}\right|\right\}^{1/\ell},
 \end{align*}
 \begin{align*}
	r_2=\max_{1\le \ell\le n}\left\{\frac{D_n}{\binom{n}{\ell}(L_{k,m-1})^{n-\ell}(L_{k,m})^{\ell}L_{k,\ell}}\left|\frac{a_{n-\ell}}{a_n}\right|\right\}^{1/\ell}.
 \end{align*}
\end{theorem}

\begin{remark}
	If $a_0=0$, then $z=0$ is already a zero of $q$ and the lower bound holds trivially with $r_1=0$; the upper bound $r_2$ does not use $a_0$ and continues to hold.
\end{remark}	
	
	\section{Preliminaries}\label{prel}
In this section, we prove some auxiliary results that will be helpful in the proof of the main theorems.
\begin{lemma}[Positivity]
	\label{lem:positivity}
	For every real $k>0$: $F_{k,n}>0$ for all integers $n\ge1$, and $L_{k,n}>0$ for all integers $n\ge0$.
\end{lemma}

\begin{proof}
	Both sequences satisfy the recurrence $x_{n+1}=kx_n+x_{n-1}$ with $k>0$. For $L$: $L_{k,0}=2>0$ and $L_{k,1}=k>0$; if $L_{k,n-1}>0$ and $L_{k,n}>0$ then $L_{k,n+1}=kL_{k,n}+L_{k,n-1}>0$, so $L_{k,n}>0$ for all $n\ge0$ by induction. For $F$: $F_{k,1}=1>0$ and $F_{k,2}=k>0$; if $F_{k,n-1}\ge0$ and $F_{k,n}>0$ then $F_{k,n+1}=kF_{k,n}+F_{k,n-1}>0$, so $F_{k,n}>0$ for all $n\ge1$ (with $F_{k,0}=0$ the sole non-positive term).
\end{proof}

\begin{cor}
	\label{cor:Dnpos}
	For every real $k>0$ and integers $m\ge1,n\ge1$, $D_n(k,m)>0$.
\end{cor}

\begin{proof}
	By Theorem \ref{thm:main},
	\[
	D_n=\sum_{\ell=1}^n\binom{n}{\ell}(L_{k,m-1})^{n-\ell}(L_{k,m})^{\ell}L_{k,\ell}.
	\]
	Each summand is a product of a positive binomial coefficient with nonnegative-integer powers of $L_{k,m-1}>0$ and $L_{k,m}>0$ (Lemma \ref{lem:positivity}, using $m-1\ge0$) and a factor $L_{k,\ell}>0$ (Lemma \ref{lem:positivity}). Hence every summand is positive, and $D_n$, a sum of $n\ge1$ such terms, is positive.
\end{proof}

\begin{lemma}[Key identity]
	\label{lem:key}
Let $\alpha,\beta$ be as in Section~\ref{sec:background}, and let $m\ge1$ be an integer. Then
	\[
	L_{k,m}-\beta L_{k,m-1}=\alpha^{m-1}(\alpha-\beta),
	\qquad
	L_{k,m}-\alpha L_{k,m-1}=-\beta^{m-1}(\alpha-\beta).
	\]
\end{lemma}

\begin{proof}
	Using $L_{k,j}=\alpha^j+\beta^j$,
	\[
	L_{k,m}-\beta L_{k,m-1}=(\alpha^m+\beta^m)-\beta(\alpha^{m-1}+\beta^{m-1})=\alpha^m-\beta\alpha^{m-1}=\alpha^{m-1}(\alpha-\beta).
	\]
	The second identity follows by exchanging $\alpha\leftrightarrow\beta$ in the same computation.
\end{proof}

\begin{lemma}[Zeros from positive weights, Dalal and Govil \cite{dalalgovil}]
	\label{lem:general}
	Let $q(z)=a_0+a_1z+\cdots+a_nz^n$ be a polynomial of degree $n\ge1$ with complex coefficients, $a_\ell\neq0$ for $0\le\ell\le n$. Suppose $B_1,\dots,B_n$ are positive reals with $\sum_{\ell=1}^nB_\ell=1$, and set
	\[
	r_1=\min_{1\le\ell\le n}\left\{B_\ell\left|\frac{a_0}{a_\ell}\right|\right\}^{1/\ell},
	\qquad
	r_2=\max_{1\le\ell\le n}\left\{\frac{1}{B_\ell}\left|\frac{a_{n-\ell}}{a_n}\right|\right\}^{1/\ell}.
	\]
	Then every zero of $q$ lies in the annulus $r_1\le|z|\le r_2$.
\end{lemma}

\begin{proof}
	This is the result of Dalal and Govil \cite{dalalgovil}; we include the short proof for completeness, adapted to the present notation. Suppose $|z|>r_2$. By definition of $r_2$, for each $\ell=1,\dots,n$,
	\[
	\left|\frac{a_{n-\ell}}{a_n}\right|\le B_\ell\, r_2^{\,\ell} < B_\ell\,|z|^{\ell},
	\]
	so $|a_{n-\ell}/a_n|\,|z|^{-\ell}<B_\ell$. Then
	\[
	|q(z)|\ge|a_n||z|^n-\sum_{j=0}^{n-1}|a_j||z|^j
	=|a_n||z|^n\left(1-\sum_{\ell=1}^n\left|\frac{a_{n-\ell}}{a_n}\right||z|^{-\ell}\right)
	>|a_n||z|^n\left(1-\sum_{\ell=1}^nB_\ell\right)=0,
	\]
	so $q(z)\neq0$. Hence every zero satisfies $|z|\le r_2$.
	
	For the lower bound, set $Q(z)=z^nq(1/z)=\sum_{\ell=0}^na_{n-\ell}z^\ell$, a degree-$n$ polynomial with nonzero coefficients (since $a_0,\dots,a_n$ are all nonzero) and leading coefficient $a_0$. Applying the bound just proved to $Q$, with the same weights $B_\ell$,
	\[
	\max_{1\le\ell\le n}\left\{\frac1{B_\ell}\left|\frac{a_\ell}{a_0}\right|\right\}^{1/\ell}
	\]
	bounds every zero $w$ of $Q$ from above; this quantity equals $1/r_1$. Since $a_0\neq0$, $0$ is not a zero of $q$, and $Q(w)=w^nq(1/w)$ shows $w$ is a zero of $Q$ iff $1/w$ is a zero of $q$. So every zero $z=1/w$ of $q$ satisfies $|z|=1/|w|\ge r_1$.
\end{proof}
	
\section{Proofs of Theorems \ref{thm:main} and \ref{thm:annulus}}

	\begin{proof}[Proof of Theorem \ref{thm:main}]
		Write $S=\sum_{\ell=0}^n\binom{n}{\ell}(L_{k,m-1})^{n-\ell}(L_{k,m})^{\ell}L_{k,\ell}$. Using $L_{k,\ell}=\alpha^{\ell}+\beta^{\ell}$,
		\[
		S=\sum_{\ell=0}^n\binom{n}{\ell}(L_{k,m-1})^{n-\ell}(L_{k,m})^{\ell}\alpha^{\ell}
		+\sum_{\ell=0}^n\binom{n}{\ell}(L_{k,m-1})^{n-\ell}(L_{k,m})^{\ell}\beta^{\ell}
		=:S_1+S_2.
		\]
		By the binomial theorem,
		\[
		S_1=(L_{k,m-1}+\alpha L_{k,m})^n,\qquad S_2=(L_{k,m-1}+\beta L_{k,m})^n.
		\]
		Since $\alpha,\beta$ are the roots of $x^2-kx-1=0$, we have $\alpha\beta=-1$. Hence, by Lemma \ref{lem:key},
		\[
		L_{k,m-1}+\alpha L_{k,m}=\alpha\bigl(L_{k,m}-\beta L_{k,m-1}\bigr)=\alpha\cdot\alpha^{m-1}(\alpha-\beta)=\alpha^{m}(\alpha-\beta),
		\]
		\[
		L_{k,m-1}+\beta L_{k,m}=\beta\bigl(L_{k,m}-\alpha L_{k,m-1}\bigr)=\beta\cdot\bigl(-\beta^{m-1}(\alpha-\beta)\bigr)=-\beta^{m}(\alpha-\beta).
		\]
		Therefore
		\[
		S_1=\alpha^{mn}(\alpha-\beta)^n,\qquad S_2=(-1)^n\beta^{mn}(\alpha-\beta)^n,
		\]
		so
		\[
		S=(\alpha-\beta)^n\bigl(\alpha^{mn}+(-1)^n\beta^{mn}\bigr).
		\]
		Since $\alpha-\beta=\sqrt{k^2+4}$, we have $(\alpha-\beta)^n=(k^2+4)^{n/2}$.
		
		If $n$ is even, then $\alpha^{mn}+(-1)^n\beta^{mn}=\alpha^{mn}+\beta^{mn}=L_{k,mn}$, and $S=(k^2+4)^{n/2}L_{k,mn}$.
		
		If $n$ is odd, then $\alpha^{mn}+(-1)^n\beta^{mn}=\alpha^{mn}-\beta^{mn}=(\alpha-\beta)F_{k,mn}$, using the Binet formula for $F_{k,mn}$, and so
		\[
		S=(\alpha-\beta)^n\cdot(\alpha-\beta)F_{k,mn}=(\alpha-\beta)^{n+1}F_{k,mn}=(k^2+4)^{(n+1)/2}F_{k,mn}.
		\]
		This establishes both cases.
	\end{proof}
	
	\begin{proof}[Proof of Theorem \ref{thm:annulus}]
		For $1\le\ell\le n$, set
		\[
		B_\ell:=\binom{n}{\ell}\frac{(L_{k,m-1})^{n-\ell}(L_{k,m})^{\ell}L_{k,\ell}}{D_n}.
		\]
		By Lemma \ref{lem:positivity} (with $m-1\ge0$) and Corollary \ref{cor:Dnpos}, $B_\ell$ is a ratio of positive quantities, so $B_\ell>0$. By Definition \ref{def:Dn} and Theorem \ref{thm:main},
		\[
		D_n=\sum_{\ell=1}^n\binom{n}{\ell}(L_{k,m-1})^{n-\ell}(L_{k,m})^{\ell}L_{k,\ell},
		\]
		so
		\[
		\sum_{\ell=1}^nB_\ell=\frac{1}{D_n}\sum_{\ell=1}^n\binom{n}{\ell}(L_{k,m-1})^{n-\ell}(L_{k,m})^{\ell}L_{k,\ell}=\frac{D_n}{D_n}=1.
		\]
		Thus $B_1,\dots,B_n$ satisfy the hypotheses of Lemma \ref{lem:general}, and substituting them into the $r_1,r_2$ of that lemma reproduces exactly the $r_1,r_2$ of Theorem \ref{thm:annulus} (since $1/B_\ell$ is precisely the reciprocal weight appearing there). The conclusion follows.
	\end{proof}

\section{Examples}
\label{sec:examples}

We illustrate Theorem \ref{thm:annulus} with one example for each parity of $n$, taking $k=1$ (so $L_{1,n}$ and $F_{1,n}$ are the classical Lucas and Fibonacci numbers) and $m=1$, and compare against Cauchy's classical bound $|z|\le1+\max_{0\le j<n}|a_j/a_n|$ recalled in the introduction.

\subsection*{Example 1 ($n=4$, even)}

Let $q(z)=1-z+2z^2+z^3+z^4$, so $(a_0,a_1,a_2,a_3,a_4)=(1,-1,2,1,1)$. With $k=1,m=1$: $L_{1,0}=2$, $L_{1,1}=1$, and
\[
D_4=(1^2+4)^{2}L_{1,4}-2(L_{1,0})^4=25\cdot7-2\cdot16=143.
\]
The weights of Theorem \ref{thm:annulus} are
\[
B_1=\frac{32}{143},\quad B_2=\frac{72}{143},\quad B_3=\frac{32}{143},\quad B_4=\frac{7}{143},
\]
which sum to $1$, giving
\[
r_1\approx0.2238,\qquad r_2\approx4.4688.
\]
The four zeros of $q$ are $-0.809\pm1.401i$ and $0.309\pm0.535i$, of moduli $1/\varphi\approx0.618$ and $\varphi\approx1.618$ respectively, where $\varphi=(1+\sqrt5)/2$ is the golden ratio; all four lie inside $[r_1,r_2]$. Cauchy's bound gives only $|z|\le3$; Theorem \ref{thm:annulus} additionally excludes $|z|<0.2238$, information Cauchy's bound cannot provide at all.

\subsection*{Example 2 ($n=5$, odd)}

Let $q(z)=2+z-z^2+z^3-z^4+z^5$, so $(a_0,\dots,a_5)=(2,1,-1,1,-1,1)$. With $k=1,m=1$ and $n$ odd, the identity now involves $F_{1,5}=5$:
\[
D_5=(1^2+4)^{3}F_{1,5}-2(L_{1,0})^5=125\cdot5-2\cdot32=561,
\]
\[
B_1=\frac{80}{561},\quad B_2=\frac{240}{561},\quad B_3=\frac{160}{561},\quad B_4=\frac{70}{561},\quad B_5=\frac{11}{561},
\]
giving
\[
r_1\approx0.2852,\qquad r_2\approx7.0125.
\]
The five zeros have moduli $0.709,\,1.253,\,1.253,\,1.340,\,1.340$ (to three decimals), all inside $[r_1,r_2]$; Cauchy's bound gives $|z|\le3$.

\begin{remark}
	In both examples $r_2$ is somewhat larger than Cauchy's bound: the annulus method is not designed to sharpen the classical upper bound in general, and we make no such claim here. Its value is structural rather than numerical: $r_1>0$ excludes a neighbourhood of the origin, which a one-sided bound such as Cauchy's cannot do.
\end{remark}

\subsection*{Comparison with the direct $k$-Fibonacci analogue}

A natural question is why the $k$-Lucas weights of Theorem \ref{thm:annulus} are worth having alongside the $k$-Fibonacci-weighted construction obtained by the same method with $L_{k,\cdot}$ replaced by $F_{k,\cdot}$ throughout (this is the direct analogue, at general $m$, of the $m=4$ bounds of D\'iaz-Barrero, Bidkham-Shashahani, and Kaur). Two points answer this.

First, at $m=1$ the $k$-Fibonacci-weighted construction is not merely weaker, it is undefined: since $F_{k,0}=0$, the normalising constant in that construction vanishes and the weights cannot be formed. The $k$-Lucas construction has no such restriction (Lemma \ref{lem:positivity}), so $m=1$, the case appearing in both examples above, is only available through Theorem \ref{thm:annulus}.

Second, where both constructions are defined, neither dominates the other. Applying both to the polynomials of Examples 1 and 2 above, now at $m=2$ (so both are defined), gives
\[
\text{Example 1: } k\text{-Lucas } [0.0102,\,97.75], \quad k\text{-Fibonacci } [0.1905,\,5.25],
\]
\[
\text{Example 2: } k\text{-Lucas } [0.0044,\,458.2], \quad k\text{-Fibonacci } [0.1818,\,11.0],
\]
where on both examples the $k$-Fibonacci-weighted bound happens to be tighter at $m=2$. This is not universal: a systematic search over 3000 random polynomials of degree $3$ to $6$, with $k\in\{0.5,1,1.5,2,3,4\}$ and $m\in\{1,2,3\}$ (1837 trials after excluding near-degenerate coefficients), found the $k$-Lucas bound gives a strictly tighter $r_1$ in $13.7\%$ of cases and a strictly tighter $r_2$ in $14.2\%$ of cases, with the $k$-Fibonacci-weighted bound tighter otherwise. The two weight families are therefore genuinely complementary rather than one subsuming the other: each is preferable on a nontrivial, non-overlapping set of polynomials, and $m=1$ is accessible only through the $k$-Lucas construction.

	
\section*{Acknowledgments}
The author is deeply grateful to his supervisor, the late Professor Florian Luca, whose foundational contributions to the study of Lucas sequences and linear recurrences shaped the direction of this work, and to whose memory this paper is dedicated. The author also thanks the Mathematics Division of Stellenbosch University for funding his PhD studies.
	
\begin{thebibliography}{99}
	
\bibitem{bidkham} 
Bidkham, M., \& Shashahani, E. (2011). An annulus for the zeros of polynomials. \textit{Appl. Math. Lett.}, \textbf{24}(2), 122--125.

\bibitem{dalalgovil}
Dalal, A., \& Govil, N. K. (2013). On region containing all the zeros of a polynomial. \textit{Appl. Math. Comput.}, \textbf{219}(17), 9609--9614.

\bibitem{diazbarrero} 
D\'iaz-Barrero, J. L. (2002). An annulus for the zeros of polynomials. \textit{J. Math. Anal. Appl.}, \textbf{273}, 349--352.
		
\bibitem{falconlucas} 
Falc\'on, S. (2011). On the $k$-Lucas numbers. \textit{Int. J. Contemp. Math. Sciences}, \textbf{6}(21), 1039--1050.

\bibitem{falconplaza} 
Falc\'on, S., \& Plaza, \'A. (2007). The $k$-Fibonacci sequence and the Pascal 2-triangle. \textit{Chaos, Solitons} \& \textit{Fractals}, \textbf{33}(1), 38--49.


\bibitem{kaur} 
Kaur, S. $k$-Fibonacci Annulus for Polynomial Zeros. Accepted by \textit{Rendiconti del Circolo Matematico di Palermo Series 2 }.

\end{thebibliography}
	
	\section*{Addresses}
	$ ^{1} $ Mathematics Division, Stellenbosch University, Stellenbosch, South Africa.
	
	Email: \url{hbatte91@gmail.com}
\end{document}
